\section{Limitations and threats to validity}
\label{sec:limits}

\subsection{No figure here is a proved minimum}

Every pair we report is a constructed upper bound. In particular, the
statement that $20$ names suffice to flatten the system of \cite{jsww1976} to
degree~$2$, and that no smaller set does, is an optimality statement
\emph{about our encoding}: the solver reports that no smaller selection exists
within the candidate catalogue we supplied, under the reduction rules we
encoded. It does not follow that no smaller flattening exists, because a
flattening may in principle use a name that is not a subexpression of the
syntactic tree, and our catalogue is built from that tree.

This is not a theoretical caveat. In an earlier line of work the reported
optimum fell repeatedly as the catalogue was enlarged---$46\to17\to15$
names---never because the search itself improved; and two of those certificates
were subsequently found not to be materializable at all (see below), so the
history cuts both ways. For the present construction we enlarged the catalogue
from $465$ to $946$ to $3469$ candidates and the optimum did not move from
$20$, but that is evidence, not proof.

\subsection{Figures we withdrew}

The development that produced the results above also produced nine defects that
we found and repaired, four of which were capable of yielding figures that
compiled, passed tests, and were wrong. We list the consequential ones because
a reader deserves to calibrate on them.

\begin{itemize}[nosep]
  \item A rewriting route in the optimizer certified sets of names that could
        not be \emph{materialized} into an actual system. This invalidated four
        published-in-progress figures at degree~$5$: $(41,5)$, $(38,5)$,
        $(36,5)$ and $(33,5)$. The frontier was re-measured in full without
        rewriting and with a new structural check; the $(44,5)$ of this paper is
        from that re-measurement.
  \item Two runs of the same sweep in one process returned $(25,13)$ and
        $(27,13)$: a wall-clock solver timeout combined with the rewriting route
        made results non-reproducible. Reproducibility across repeated runs is
        now itself checked.
  \item A list of proved nonnegativity facts was matched to expressions by
        string comparison; after a reparametrization the strings no longer
        matched and the facts were silently lost, making a whole corner appear
        unsatisfiable. Facts are now attached to systems structurally.
  \item Three statement defects in the Lean generator, described in
        \S\ref{sec:lean}, two of which produced files that compiled.
\end{itemize}

The safeguards these prompted---generated rather than transcribed formal
statements, statement tests that check the text of each proof file against the
published equations, an axiom audit in the build, and reproducibility
checks---are what we would ask a reader to weigh, rather than the absence of
defects.

\subsection{Upper bounds on degrees in \texorpdfstring{\S\ref{sec:sec3}}{Section 6}}

The degrees of the relation-combining curve are computed by traversing the
expression tree without expanding, because expansion is infeasible: substituting
the definitions of \S\ref{sec:sec3} in cascade doubles the degree at each step
and our computer algebra system does not terminate. They are therefore upper
bounds, exact only if no leading terms cancel. Four independent published
figures are reproduced exactly by the same computation, which is our evidence
that no cancellation occurs at those points; cancellation elsewhere would
\emph{lower} the reported degrees.

\subsection{The $(51,5)$ rests on the same kind of instrument}

The exact minimum of $25$ names for the method of \cite{davis1973}, which
supports our strongest claim about the literature (\S\ref{sec:correc}), comes
from the same class of tool as the $20$: an SMT optimizer, over our
formalization of the method, over our transcription of the system. What differs
is that the method's admissible names are \emph{exactly} the degree-$2$
monomials, so the candidate set is determined rather than chosen---which is why
we call the minimum exact. But the lower bound is still a solver's, and this
project has a documented case in which a certified lower bound from a related
encoding sat above a flattening we later constructed by hand. We ran the
computation on two syntactic forms of the system and obtained $25$ with a
matching bound both times.

\subsection{Only the variable count is formalized}

For each pair $(\nu,\delta)$ we report, the formal development establishes the
equisatisfiability, and hence $\nu$. The degree $\delta$ is computed by our
tooling from the transformed system and checked against
Proposition~\ref{prop:gen}, but it is not a Lean theorem. A reader who wants
$\delta$ machine-checked does not have it here.

\subsection{The four points of \texorpdfstring{\S\ref{sec:sec3}}{Section 6} are announced, not exhibited}

We apply our own criterion to ourselves. The pairs $(16,369)$, $(15,801)$,
$(14,1777)$ and $(13,3905)$ are computed, not written down: no polynomial is
displayed, nothing is formalized, and the degrees are upper bounds. They belong
in the same category as the announcements this paper audits, and should be cited
that way.

\subsection{A domination we do not claim}

The generator of \cite{matiyasevich1981} is quoted as having degree
\emph{greater than} $6000$. That is a lower bound, so it settles nothing about
whether its degree is below or above the $13697$ of the $12$-variable
polynomial. We therefore do not assert that it dominates any of our points on
both axes---only that it has fewer variables. The exact degree is recoverable
from the Mizar formalization \cite{pak2022ten}; extracting it would close the
comparison, and we have not done so.

\subsection{What is cited and not reproved}

We do not reprove that the system of \cite{jsww1976} represents the primes. That
theorem is theirs and we cite it. What we verify is that our transformations of
their system preserve satisfiability, and that our transcription of their system
is faithful, equation by equation. A reader who doubts their theorem should
doubt our figures equally.

Finally, the constructions of \S\ref{sec:sec3} rest on our transcription of
their Theorem~3.9 into polynomial form, which introduces eight unknowns (six
square roots, one quotient, one slack). We showed in \S\ref{sec:sec3} that this
conversion carries no additional hypothesis, but it is a conversion, and its
faithfulness is checked by the same statement tests rather than proved.
